\chapter{Business Insights and Decision-Making}\label{chap:business_decisions}

\section{From Customer Segments to Business Actions}
The three dashboard pages support different but connected questions. The executive page shows the distribution of customer segments over time, the demographic page helps the team understand whom each segment contains, and the RFM page shows how recently and frequently customers buy and how much they spend. These results can be used to prioritize marketing resources rather than sending the same campaign to the entire customer base.

The strategies below have been retained from the team's dashboard analysis. They are \emph{proposed business actions}, not campaigns executed or experimentally proven by this project. Reported numerical values refer to the dashboard's 30 June 2014 analytical snapshot and should be reconciled against the final refresh before presentation.

\section{Actionable Business Recommendations}
\label{sec:bi_recommendations}
Based on the dashboard insights, we recommend the following practical strategies for CompanyX:

\subsection{1. Champions: VIP Loyalty Program}
\label{subsec:rec_champions}
\begin{itemize}
    \item \textbf{Goal:} Keep our best customers happy and encourage them to recommend the brand to others.
    \item \textbf{Action:} Create a VIP club that gives early access to new bicycle models, invitations to special cycling events, and referral discounts when they introduce friends.
\end{itemize}

\subsection{2. High-Value / Cooling Off: Urgent Retention Campaigns}
\label{subsec:rec_high_value}
\begin{itemize}
    \item \textbf{Goal:} Protect a segment representing over 71\% of the monetary amount in the final reporting snapshot.
    \item \textbf{Action:} Set up an automated email at day 120 after their purchase offering a free annual bicycle check-up at local partner shops, or a \$150 discount voucher if they upgrade components within 30 days.
\end{itemize}

\subsection{3. Recent / Promising: Cross-Selling and Welcome Series}
\label{subsec:rec_recent}
\begin{itemize}
    \item \textbf{Goal:} Encourage new buyers to make a second purchase quickly.
    \item \textbf{Action:} Send a 30-day email series: Day 3 with bike assembly and safety tips, Day 14 with cleaning advice, and Day 21 offering a 15\% discount on essential accessories like helmets, gloves, and lights.
\end{itemize}

\subsection{4. Low-Spend / Occasional: Low-Cost Automated Marketing}
\label{subsec:rec_low_spend}
\begin{itemize}
    \item \textbf{Goal:} Increase purchase frequency without spending too much marketing budget.
    \item \textbf{Action:} Avoid expensive phone calls or printed mail. Instead, send automated emails and seasonal clearance announcements for cheap consumables like water bottles and spare tubes.
\end{itemize}

\subsection{5. Lost / Dormant: Cost Reduction}
\label{subsec:rec_lost}
\begin{itemize}
    \item \textbf{Goal:} Avoid costly broad advertising to a low-engagement segment (whose average snapshot recency exceeds 600 days).
    \item \textbf{Action:} Remove this group from paid Facebook and Google ads. Only send them one or two broad holiday sale emails per year (e.g., Black Friday or New Year).
\end{itemize}


\section{Business Decision-Making Use Case}
\label{sec:bi_usecase}
Consider a marketing team with a limited campaign budget. The dashboard shows that the \emph{High-Value / Cooling Off} segment contains 7,357 customers and represents a large share of the June 2014 snapshot's monetary value. The team can use the segment slicer to inspect the associated demographics, then prioritize a time-limited maintenance or accessory follow-up campaign for this group. A second, low-cost welcome campaign can target \emph{Recent / Promising} buyers, whose small initial spend and low recency suggest an opportunity to encourage a further purchase.

This is a \emph{decision-support use case}, not evidence that either campaign was executed or increased revenue. An eventual campaign evaluation would compare actual response, incremental sales, and campaign costs against a suitable baseline. The dashboard therefore demonstrates how the segmentation provides a targetable business action while keeping the distinction between descriptive analysis and realized business impact.

\paragraph{Interpretation note.} The dashboard's figures reflect the current saved report and reference-date filters. The ML chapter contains an illustrative 18,484-customer model extract, whereas this dashboard reports 18,485 customers. The team should reconcile that one-record discrepancy during the final dashboard refresh before asserting identical populations across the two exports.
